\section{Introduction}
Cardiovascular diseases remain one of the leading causes of mortality worldwide.
Early detection and risk assessment play a crucial role in preventing severe outcomes and improving
patients' quality of life. In recent years, data-driven methods and machine learning techniques have
become increasingly important tools for analyzing medical data and supporting clinical decision making.

In this project, we apply data mining techniques to the well-known Heart Disease dataset from the
UCI Machine Learning Repository. The main objective is to investigate how patients' characteristics,
such as age, sex, chest pain type, cholesterol level, resting blood pressure, and maximum heart rate,
can be used to predict the presence of heart disease.

We focus on the K-Nearest Neighbors (KNN) classification algorithm, which is a simple yet powerful
distance-based method. The project includes several main steps: data preprocessing, exploratory data analysis (EDA),
implementation of the KNN classifier for mixed-type attributes, and performance evaluation
using standard classification metrics. The results of this study illustrate the potential of
KNN for medical classification tasks and highlight the importance of proper preprocessing
and feature analysis in data mining workflows.